\section{Localizing the Candidate Bottleneck}
\label{sec:localization}

We use a sequence of interventions that progressively hold more of the pipeline fixed. The goal is not to turn one selected positive case into a general claim; it is to determine which layer remains plausible after simpler explanations fail.

\subsection{Evidence-source effects are small and heterogeneous at scale}
Our closed DeepSeek V2 study contains 48 matched pairs, 96 learned states, and 1,728 held-out task evaluations. The two primary states differ in the evidence projection exposed to the same frozen updater. Across pairs, the mean MRW$-$WIN-C utility difference is $+0.0231$, the median is $+0.0139$, and the standard deviation is $0.0780$. The frozen one-sided sign-flip test gives $p=0.1719$, and the bootstrap confidence interval is $[-0.0185,0.0660]$.

This study therefore does not establish rejected-witness superiority. Individual streams span positive, null, and negative effects, including values from $-0.125$ to $+0.181$. The appropriate inference is heterogeneity, not a failed attempt that can be repaired by silently adding samples or widening a threshold.

\subsection{Selecting a more ``diagnostic'' failure does not solve the problem}
Post-hoc development analysis suggested that progress-related features correlated with pairwise effects, motivating a single-stream witness-selector pilot on \texttt{e1-tsr-00}. This stream was chosen after outcome inspection and is therefore mechanism development only. We froze four evidence arms: WIN-C, First-Fail, Progress-Fail, and Progress-Contrast. On 18 held-out tasks they score 13/18, 17/18, 14/18, and 14/18 respectively.

Although the historical First-Fail state is strong, the preregistered purpose of this pilot was to test whether a progress-based selector provides a stable repair rule. It does not. The frozen verdict is \texttt{S1\_SIGNAL\_FAIL\_STOP\_NO\_S2}. We therefore reject the simple story that ``choose a later or more diagnostic failure'' explains the earlier heterogeneity.

\subsection{The strong persistent state remains useful when frozen}
The next intervention holds the persistent state bytes fixed and re-runs only downstream evaluation. The historical strong First-Fail state is compared with the corresponding WIN-C state in two fresh complete actor replicates on the same 18-task development panel. Replicate 1 gives 15/18 versus 14/18; replicate 2 gives 16/18 versus 12/18.

These remeasurements do not make the selected stream representative, and they do not prove that First-Fail evidence caused the state. They do weaken a simpler explanation: the original 17/18 First-Fail result was not merely one lucky downstream actor realization attached to an otherwise inert state. The frozen artifact itself carries behaviorally relevant information.

\subsection{Byte-identical evidence fails to regenerate the same advantage}
We then reconstruct the exact S1 First-Fail evidence rendered to the updater. All 16 rendered packet hashes match the historical inputs. From this byte-identical evidence we create two new free-form updater realizations, then evaluate each fresh state against a fresh WIN-C realization on the same 18 tasks. Replicate 1 is 15/18 versus 15/18; replicate 2 is 11/18 versus 12/18. The mean fresh contrast is $-0.0278$.

The contrast with frozen-state remeasurement is the key localization result. Holding the strong state fixed preserves a directional advantage; holding the evidence fixed but regenerating the state does not. Because the fresh updater states and their downstream evaluations both vary, this is not a full variance-component decomposition. It is nevertheless incompatible with the claim that evidence availability alone explains the strong historical state.

\begin{table}[t]
\centering
\caption{Evidence ladder. The independent unit and scientific interpretation change across stages; no row is promoted beyond its design.}
\label{tab:evidence-ladder}
\small
\setlength{\tabcolsep}{3pt}
\begin{tabular}{p{0.18\linewidth}p{0.18\linewidth}p{0.22\linewidth}p{0.27\linewidth}}
\toprule
Stage & Held fixed & Observation & Supported interpretation \\
\midrule
V2 48-pair study & actor/updater configuration, task design & mean MRW$-$WIN-C $=+0.0231$, CI crosses 0 & evidence-source effect is heterogeneous \\
S1 selector pilot & selected stream, updater path & 13/18, 17/18, 14/18, 14/18 & simple progress selector fails \\
Frozen-state remeasure & persistent state bytes & FF: 15/18, 16/18; WIN-C: 14/18, 12/18 & strong state has reproducible local behavioral value \\
Exact-evidence replay & learner-visible evidence bytes & FF: 15/18, 11/18; WIN-C: 15/18, 12/18 & same evidence does not reliably regenerate strong state \\
\bottomrule
\end{tabular}
\end{table}

\subsection{What is and is not localized}
The current evidence supports the following bounded statement:

\begin{quote}
In one controlled development case, identical trajectory evidence did not reliably regenerate the same useful persistent skill state, while a previously generated strong state retained downstream utility when frozen. Persistent-state generation is therefore a candidate bottleneck that cannot be reduced to evidence availability alone.
\end{quote}

It does \emph{not} establish that rejected failures are generally better than winners, that free-form updater variance dominates actor variance in the population, or that deterministic compilation is already superior. These stronger claims require prospective interventions rather than retrospective interpretation of the selected case.
